\documentclass[../../main.tex]{subfiles}

\begin{document}

\chapter{Mode Pengalamatan}
\label{chap:mode-pengalamatan}

% =========================================================
% OUTCOME PEMBELAJARAN KHUSUS BAB
% =========================================================

\begin{subcpmkbox}
\SubCPMK{2.2} --- Menganalisis mode pengalamatan (\eng{immediate},
\eng{direct}, \eng{register indirect}, \eng{indexed},
\eng{base-relative}, dan \eng{based-indexed}) serta menentukan
segmen \eng{default} yang digunakan setiap mode.

\vspace{0.4em}
\textbf{CPMK yang didukung:} \CPMK{02} \quad\textbullet\quad
\textbf{CPL yang didukung:} \CPL{02}
\end{subcpmkbox}

\begin{tujuanbox}
Setelah mempelajari bab ini, mahasiswa mampu:
\begin{enumerate}
    \item menjelaskan peran mode pengalamatan dalam menentukan letak
          operand;
    \item mengenali mode pengalamatan hanya dari penulisan
          operandnya;
    \item menghitung alamat efektif (\eng{effective address}) dan
          alamat fisik suatu operand;
    \item menentukan segmen \eng{default} setiap mode serta menuliskan
          \eng{segment override prefix} yang tepat;
    \item memilih mode yang paling efisien untuk mengakses elemen
          array dan medan \eng{record};
    \item menuliskan satu operasi yang sama dalam beberapa mode
          berbeda.
\end{enumerate}
\end{tujuanbox}

\begin{materibox}
\begin{enumerate}
    \item Konsep operand sumber dan tujuan serta alasan mode
          pengalamatan diperlukan.
    \item Mode \eng{immediate}, \eng{register}, dan \eng{direct}
          (memori langsung).
    \item Mode \eng{register indirect}: \code{[BX]}, \code{[SI]},
          \code{[DI]}, dan \code{[BP]}.
    \item Mode \eng{indexed}, \eng{base-relative}, dan
          \eng{based-indexed}.
    \item Kombinasi register yang legal dan yang ilegal.
    \item Segmen \eng{default} dan \eng{segment override prefix}
          (\code{CS:}, \code{DS:}, \code{SS:}, \code{ES:}).
    \item Perbandingan efisiensi: jumlah byte dan siklus relatif.
    \item Studi kasus akses array satu dimensi dan \eng{record}.
\end{enumerate}
\end{materibox}

\subfile{section-01}
\subfile{section-02}
\subfile{section-03}

\section{Contoh Program}
\label{sec:contoh-bab8}

Dua program berikut dirakit dan dijalankan dengan Emu8086 atau
TASM/TLINK. Keduanya memperlihatkan bahwa mode pengalamatan langsung
menentukan bentuk kode dan alamat yang dibaca prosesor.

\subsection{Contoh 1: Membandingkan Tiga Mode pada Array}

Program \file{mode-array.asm} menjumlahkan lima elemen array yang
sama dengan tiga mode pengalamatan berbeda, lalu membandingkan ketiga
hasilnya. Kode lengkapnya tampak pada Kode~\ref{lst:mode-array}.

\lstinputlisting[
    language=MASM,
    caption={Program \file{mode-array.asm}: tiga mode pengalamatan
    untuk satu array.},
    label={lst:mode-array}
]{\CodeDir/mode-array.asm}

Cara kerja program dijelaskan sebagai berikut.

\begin{itemize}
    \item \textbf{Cara 1 --- mode direct.} Elemen dibaca dengan offset
          tetap, \lstasm|data1+0| sampai \lstasm|data1+8|. Mudah
          dibaca, tetapi tidak praktis untuk array panjang.
    \item \textbf{Cara 2 --- mode indexed.} \reg{SI} menjadi indeks
          yang bertambah dua setiap putaran; \lstasm|data1[SI]|
          setara dengan \lstasm|[data1+SI]| pada segmen \reg{DS}.
    \item \textbf{Cara 3 --- mode base-relative.} \reg{BX} diisi
          \lstasm|LEA BX, data1| sebagai basis, lalu elemen dibaca
          dengan \lstasm|[BX+SI]|; basis dan indeks terpisah sehingga
          mudah dipakai ulang.
    \item \textbf{Pembandingan hasil.} \asm{CMP} dan \asm{JNE}
          membandingkan \code{jumlah1} dengan \code{jumlah2}, lalu
          \code{jumlah2} dengan \code{jumlah3}. Struktur kontrol ini
          dibahas pada Bab~\ref{chap:kontrol} dan operasi
          aritmatikanya pada Bab~\ref{chap:aritmatika}.
\end{itemize}

Ketiga cara menghasilkan jumlah yang identik, yaitu
$10 + 20 + 30 + 40 + 50 = 150$ atau \hex{96}.

\begin{session}[caption={Keluaran program \file{mode-array.asm}.}, label={lst:keluaran-mode-array}]
Perbandingan tiga mode pengalamatan
Ketiga hasil identik.
\end{session}

\begin{infobox}
Bila hasil ketiga cara berbeda, kesalahan yang paling sering terjadi
adalah keliru memakai \reg{BP} sebagai basis. Bentuk \lstasm|[BP+SI]|
membaca segmen \reg{SS}, bukan \reg{DS}, sehingga data yang diambil
adalah isi stack, bukan isi array.
\end{infobox}

\subsection{Contoh 2: Membaca Parameter dengan \code{[BP+disp]}}

Program \file{param-bp.asm} memperlihatkan pemakaian mode
base-relative dengan register \reg{BP} untuk membaca dua nilai yang
diletakkan pada stack. Kode lengkapnya tampak pada
Kode~\ref{lst:param-bp}.

\lstinputlisting[
    language=MASM,
    caption={Program \file{param-bp.asm}: membaca parameter melalui
    \code{[BP+disp]}.},
    label={lst:param-bp}
]{\CodeDir/param-bp.asm}

Poin penting program ini adalah sebagai berikut.

\begin{enumerate}
    \item Nilai 2 dan 1 diletakkan pada stack melalui \asm{PUSH}.
          Prosesor 8086 tidak memiliki \lstasm|PUSH 2| untuk nilai
          segera, sehingga nilainya dimuat ke \reg{AX} terlebih
          dahulu.
    \item Setelah \lstasm|MOV BP, SP|, \reg{BP} menunjuk puncak
          stack. \lstasm|MOV AX, [BP+0]| mengambil parameter pertama
          dan \lstasm|MOV BX, [BP+2]| mengambil parameter kedua.
    \item Keduanya memakai segmen \reg{SS} secara bawaan. Bila
          penulisannya diubah menjadi \lstasm|[BX+2]| tanpa
          menyiapkan \reg{BX}, program akan membaca lokasi yang
          berbeda pada segmen \reg{DS}.
    \item Hasil penjumlahan bernilai 3 dan dicetak sebagai karakter
          melalui layanan DOS fungsi \hex{02}; konversi bilangan
          dibahas pada Bab~\ref{chap:aritmatika}.
    \item Dua instruksi \asm{POP} mengembalikan \reg{SP} ke posisi
          semula, kebiasaan yang akan dipakai secara baku pada
          Bab~\ref{chap:prosedur}.
\end{enumerate}

\begin{session}[caption={Keluaran program \file{param-bp.asm}.}, label={lst:keluaran-param-bp}]
1 + 2 = 3
\end{session}

\section{Praktikum}
\label{sec:praktikum-bab8}

\begin{praktikumbox}
\textbf{Judul:} Penelusuran mode pengalamatan dan verifikasi alamat
efektif dengan \eng{debugger}.\\[0.3em]
\textbf{Sub-CPMK:} 2.2 \quad\textbullet\quad
\textbf{Alokasi waktu:} 100 menit \quad\textbullet\quad
\textbf{Bentuk:} individu
\end{praktikumbox}

\subsection{Tujuan Praktikum}

Setelah menyelesaikan praktikum ini, mahasiswa mampu menelusuri
eksekusi program satu instruksi demi satu, menghitung alamat efektif
setiap operand memori, dan memverifikasinya melalui jendela memori.

\subsection{Alat dan Bahan}

\begin{itemize}
    \item Komputer dengan sistem operasi Windows.
    \item Emu8086 (emulator dengan tampilan register dan memori).
    \item DOSBox beserta \file{DEBUG} untuk penelusuran berbasis
          teks, serta TASM dan TLINK untuk merakit program.
    \item Berkas \file{mode-array.asm} dan \file{param-bp.asm} dari
          \file{code/asm/}.
    \item Lembar kerja pengamatan alamat efektif.
\end{itemize}

\subsection{Langkah Kerja}

\begin{enumerate}
    \item \textbf{Menyiapkan program.} Rakit kedua berkas contoh
          dengan TASM dan TLINK, lalu pastikan keluarannya sesuai
          Kode~\ref{lst:keluaran-mode-array} dan
          Kode~\ref{lst:keluaran-param-bp}.

    \item \textbf{Menelusuri dengan Emu8086.}
    \begin{enumerate}
        \item Buka \file{mode-array.asm}, tekan \file{Compile and
              Emulate}, lalu aktifkan tombol \file{Single Step}.
        \item Catat nilai \reg{DS}, \reg{BX}, \reg{SI}, dan \reg{BP}
              sebelum setiap instruksi memori dieksekusi, lalu hitung
              alamat efektifnya secara manual.
        \item Bandingkan hasil hitungan Anda dengan alamat yang
              ditampilkan emulator.
    \end{enumerate}

    \item \textbf{Menelusuri dengan \file{DEBUG}.}
    \begin{enumerate}
        \item Jalankan DOSBox, lalu ketik
              \file{debug mode-array.exe}.
        \item Tekan \file{u} untuk melihat \eng{disassembly} dan
              \file{r} untuk melihat isi register.
        \item Tekan \file{t} beberapa kali dan catat perubahan
              register setiap instruksi.
        \item Tampilkan \file{d ds:0} dan \file{d ss:0}, lalu
              verifikasi alamat yang ditunjuk benar-benar berisi nilai
              array, bukan nol atau karakter sampah.
    \end{enumerate}

    \item \textbf{Membandingkan \code{[BX]} dan \code{[BP]}.}
    \begin{enumerate}
        \item Pada \file{param-bp.asm}, jalankan sampai
              \lstasm|MOV BP, SP|, lalu catat nilai \reg{SS} dan
              \reg{SP}.
        \item Ubah \lstasm|MOV AX, [BP+0]| menjadi
              \lstasm|MOV AX, [BX+0]| dengan \reg{BX} bernilai nol.
        \item Rakit ulang dan jelaskan mengapa keluarannya berubah,
              dengan mengaitkannya pada
              Tabel~\ref{tab:segmen-default}.
    \end{enumerate}

    \item \textbf{Mengisi lembar pengamatan.} Lengkapi
          Tabel~\ref{tab:pengamatan-ea} berdasarkan hasil penelusuran,
          lalu kerjakan Kuis 2 pada akhir praktikum.
\end{enumerate}

\begin{table}[H]
    \centering
    \caption{Lembar pengamatan alamat efektif yang harus diisi
    mahasiswa.}
    \label{tab:pengamatan-ea}
    \small
    \begin{tabular}{@{}p{0.7cm}p{3.7cm}p{2.7cm}p{1.7cm}p{2.2cm}p{1.8cm}@{}}
        \toprule
        \textbf{No} & \textbf{Instruksi} & \textbf{EA terhitung} &
        \textbf{Segmen} & \textbf{Isi memori} &
        \textbf{Sesuai} \\
        \midrule
        1 & \lstasm|MOV AX, data1+0| & \code{offset data1} & \reg{DS}
        & \hex{000A} (10) & Ya \\
        \addlinespace
        2 & \lstasm|MOV AX, data1[SI]|, \reg{SI} = 2 &
        \dots & \dots & \dots & \dots \\
        \addlinespace
        3 & \lstasm|MOV AX, [BX+SI]|, \reg{SI} = 4 &
        \dots & \dots & \dots & \dots \\
        \addlinespace
        4 & \lstasm|MOV AX, [BP+0]| pada \file{param-bp.asm} &
        \dots & \dots & \dots & \dots \\
        \addlinespace
        5 & \lstasm|MOV AX, [BP+2]| pada \file{param-bp.asm} &
        \dots & \dots & \dots & \dots \\
        \bottomrule
    \end{tabular}
\end{table}

\subsection{Hasil yang Dikumpulkan}

\begin{enumerate}
    \item Tabel~\ref{tab:pengamatan-ea} yang telah dilengkapi.
    \item Tangkapan layar jendela register dan memori Emu8086 untuk
          minimal tiga instruksi berbeda.
    \item Tangkapan layar sesi \file{DEBUG} dengan perintah
          \file{u}, \file{t}, dan \file{d}, serta lembar jawaban
          Kuis 2.
\end{enumerate}

\subsection{Kuis 2}

Kuis 2 dilaksanakan pada 15 menit terakhir praktikum secara individu
tanpa membuka catatan.

\begin{enumerate}
    \item Tentukan mode pengalamatan dari lima penulisan berikut:
          \lstasm|MOV AX, 20|, \lstasm|MOV AX, BX|,
          \lstasm|MOV AX, [BX]|, \lstasm|MOV AX, [SI+4]|, dan
          \lstasm|MOV AX, [BX+DI+2]|.

    \item Diketahui \reg{DS} = \hex{2000}, \reg{BX} = \hex{0100},
          \reg{SI} = \hex{0004}, dan \reg{SS} = \hex{3000}.
          Hitung alamat efektif dan alamat fisik untuk
          \lstasm|MOV AX, [BX+SI+8]| serta
          \lstasm|MOV AX, [BP]| dengan \reg{BP} = \hex{0010}.

    \item Program membaca variabel segmen data dengan
          \lstasm|MOV AX, [BP+2]| dan hasilnya salah. Jelaskan
          penyebabnya dan tuliskan dua cara memperbaikinya tanpa
          mengubah nilai perpindahannya.
\end{enumerate}

\begin{asesmenbox}
\textbf{Asesmen praktikum 2 dan Kuis 2 (bobot 10\% dari komponen
praktikum).}

\begin{tabular}{@{}p{5.9cm}cp{4.9cm}@{}}
    \toprule
    \textbf{Indikator} & \textbf{Bobot} & \textbf{Kriteria} \\
    \midrule
    Ketepatan perhitungan alamat efektif & 25\% &
    Seluruh baris lembar pengamatan terisi benar \\
    \addlinespace
    Verifikasi melalui jendela memori & 20\% &
    Alamat dan isi memori yang dicatat sesuai dengan
    \eng{debugger} \\
    \addlinespace
    Analisis perbedaan \code{[BX]} dan \code{[BP]} & 15\% &
    Penjelasan menyebut peran segmen \reg{DS} dan \reg{SS} \\
    \addlinespace
    Kuis 2 & 30\% &
    Minimal tiga dari empat bagian soal terjawab benar \\
    \addlinespace
    Kerapian dokumentasi praktikum & 10\% &
    Tangkapan layar lengkap dan langkah kerja dapat diulang \\
    \bottomrule
\end{tabular}
\end{asesmenbox}

\section{Latihan}
\label{sec:latihan-bab8}

\begin{enumerate}
    \item \textbf{MOV.} Jelaskan aturan operand pada instruksi
          \asm{MOV}. Sebutkan tiga bentuk pemindahan yang
          diperbolehkan dan dua bentuk yang tidak diperbolehkan.

    \item \textbf{Perbaikan program.} Instruksi
          \lstasm|MOV [BX], [SI]| tidak dapat dirakit. Tuliskan
          penggantinya dengan register perantara, lalu jelaskan
          mengapa cara itu diperlukan.

    \item \textbf{LEA.} Jelaskan perbedaan
          \lstasm|LEA SI, nilai| dan \lstasm|MOV SI, nilai|.
          Kapan \asm{LEA} lebih tepat dipakai?

    \item \textbf{Far pointer.} Diketahui data
          \code{ptr DD 12345678h}.
          \begin{enumerate}
              \item Tuliskan empat byte yang tersimpan di memori.
              \item Tentukan isi \reg{SI} dan \reg{DS} setelah
                    \lstasm|LDS SI, ptr|.
              \item Tentukan isi \reg{DI} dan \reg{ES} setelah
                    \lstasm|LES DI, ptr|.
              \item Jelaskan mengapa \asm{LDS} harus dipakai dengan
                    hati-hati.
          \end{enumerate}

    \item \textbf{XCHG dan CMP.} Diketahui \reg{AX} = \hex{1234}
          dan \reg{BX} = \hex{5678}.
          \begin{enumerate}
              \item Tentukan isi \reg{AX} dan \reg{BX} setelah
                    \lstasm|XCHG AX, BX|.
              \item Tentukan \eng{flag} yang berubah setelah
                    \lstasm|CMP AX, BX|.
              \item Jelaskan apakah \asm{CMP} mengubah isi kedua
                    operand.
          \end{enumerate}

    \item \textbf{XLAT.} Diketahui tabel
          \code{DB '0','1','2','3','4','5','6','7','8','9'},
          \reg{BX} menunjuk awal tabel, dan \reg{AL} = 7.
          \begin{enumerate}
              \item Tentukan isi \reg{AL} setelah \asm{XLATB}.
              \item Jelaskan peran \reg{BX} dan \reg{AL} pada
                    instruksi tersebut.
              \item Sebutkan satu kegunaan \asm{XLAT} dalam program.
          \end{enumerate}

    \item \textbf{CBW dan CWD.} Tentukan hasil setiap instruksi.
          \begin{enumerate}
              \item \reg{AL} = \hex{7F}, lalu \asm{CBW}
              \item \reg{AL} = \hex{80}, lalu \asm{CBW}
              \item \reg{AX} = \hex{7FFF}, lalu \asm{CWD}
              \item \reg{AX} = \hex{8000}, lalu \asm{CWD}
          \end{enumerate}
          Jelaskan mengapa hasilnya berbeda untuk bilangan positif
          dan negatif.

    \item \textbf{Penulisan program.} Tulislah program yang:
          \begin{enumerate}
              \item mendefinisikan dua variabel word;
              \item menyalin keduanya ke register;
              \item menukar isi kedua register dengan \asm{XCHG};
              \item membandingkan keduanya dengan \asm{CMP};
              \item menyimpan hasil pertukaran ke variabel.
          \end{enumerate}
\end{enumerate}

\begin{refleksibox}
\vspace{1.1em}
Jawab pertanyaan berikut pada jurnal praktikum Anda.
\begin{enumerate}
    \item Instruksi mana yang paling sulit dipahami: \asm{LEA},
          \asm{LDS/LES}, \asm{XLAT}, atau \asm{CBW/CWD}? Mengapa?
    \item Setelah menelusuri program, apakah Anda dapat membedakan
          dengan jelas antara alamat dan isi memori? Jelaskan.
    \item Konsep apa yang perlu Anda ulangi sebelum mempelajari
          interupsi BIOS/DOS pada bab berikutnya?
\end{enumerate}
\end{refleksibox}

\section{Rangkuman}
\label{sec:rangkuman-bab8}

Pada bab ini telah dibahas:

\begin{itemize}
    \item \term{Mode pengalamatan} adalah cara prosesor menyatakan
          letak operand. Ia menjadi jembatan antara penulisan
          simbolik programmer dan byte \eng{ModR/M} yang dihasilkan
          assembler.

    \item Tujuh mode pengalamatan 8086 adalah \eng{immediate},
          \eng{register}, \eng{direct}, \eng{register indirect},
          \eng{indexed}, \eng{base-relative}, dan
          \eng{based-indexed}.

    \item Mode \eng{immediate} dan \eng{register} tidak mengakses
          memori sehingga paling cepat. Mode \eng{direct} memakai
          offset tetap, sedangkan empat mode sisanya membentuk alamat
          dari register penunjuk.

    \item \term{Alamat efektif} dihitung sebagai jumlah basis,
          indeks, dan perpindahan: \code{EA = base + index + disp}.
          Alamat fisik kemudian dibentuk dengan menambahkan basis
          segmen, yaitu \code{(segmen x 16) + EA}.

    \item Hanya \reg{BX}, \reg{BP}, \reg{SI}, dan \reg{DI} yang boleh
          menjadi penunjuk alamat. Kombinasi seperti \code{[AX]},
          \code{[BX+BP]}, \code{[SI+DI]}, dan \code{[BX+SI+DI]}
          ditolak karena perangkat keras tidak menyediakan
          \eng{encoding}-nya.

    \item Segmen \eng{default} bergantung pada register penunjuk:
          kelompok \reg{DS} untuk \reg{BX}, \reg{SI}, dan \reg{DI};
          kelompok \reg{SS} untuk \reg{BP}. \eng{Segment override
          prefix} seperti \code{ES:} mengubah segmen yang dipakai
          tanpa mengubah nilai offsetnya.

    \item Efisiensi menurun seiring bertambahnya komponen alamat.
          Mode \eng{register} memerlukan dua siklus, sedangkan mode
          \eng{based-indexed} memerlukan sekitar sembilan belas
          siklus untuk instruksi yang sama panjangnya.

    \item Mode \eng{indexed} dipakai untuk menelusuri elemen array,
          sedangkan mode \eng{base-relative} dipakai untuk membaca
          medan \eng{record} dan parameter prosedur melalui
          \code{[BP+disp]}.
\end{itemize}

\begin{checklistbox}
Tandai setiap butir yang sudah Anda kuasai.

\begin{itemize}
    \item[\checkbox] Saya dapat menyebutkan tujuh mode pengalamatan
          8086 beserta contoh penulisannya.
    \item[\checkbox] Saya dapat menentukan mode pengalamatan hanya
          dari penulisan operandnya.
    \item[\checkbox] Saya dapat menghitung alamat efektif dan alamat
          fisik suatu operand.
    \item[\checkbox] Saya dapat menjelaskan mengapa \code{[AX]} dan
          \code{[SI+DI]} tidak diperbolehkan.
    \item[\checkbox] Saya dapat menentukan segmen \eng{default} dan
          menuliskan \eng{segment override prefix} dengan benar.
    \item[\checkbox] Saya dapat memilih mode yang paling efisien untuk
          mengakses elemen array dan medan \eng{record}.
    \item[\checkbox] Saya berhasil menelusuri alamat efektif dengan
          \eng{debugger} dan memverifikasinya pada jendela memori.
\end{itemize}

\vspace{0.4em}
Bila masih ada butir yang belum bertanda, ulangi bagian yang
bersesuaian sebelum melanjutkan ke Bab~\ref{chap:translasi}.
\end{checklistbox}

\begin{infobox}
\textbf{Rujukan bab.} Materi pada bab ini diperkuat oleh
\cite{intel1979} untuk daftar \eng{encoding} dan waktu siklus,
\cite{kader2014} untuk pengelompokan mode pengalamatan,
\cite{evans2020} untuk aturan segmen \eng{default} dan
\eng{override}, serta \cite{uva2024} untuk latihan penelusuran alamat
efektif. Untuk pembahasan segmen secara lebih mendalam, lihat
Bab~\ref{chap:segmentasi}.
\end{infobox}

\section{Evaluasi}
\label{sec:evaluasi-bab8}

\subsection{Soal Pilihan Ganda}

\begin{enumerate}
    \item Mode pengalamatan yang dipakai oleh operand
          \lstasm|MOV AX, [BX+SI+4]| adalah \dots
          \begin{enumerate}
              \item \eng{indexed}
              \item \eng{base-relative}
              \item \eng{based-indexed}
              \item \eng{register indirect}
          \end{enumerate}

    \item Segmen \eng{default} untuk operand \lstasm|MOV AX, [BP+2]|
          adalah \dots
          \begin{enumerate}
              \item \reg{CS} \item \reg{DS}
              \item \reg{ES} \item \reg{SS}
          \end{enumerate}

    \item Penulisan operand berikut yang ditolak assembler karena
          tidak memiliki \eng{encoding} 8086 adalah \dots
          \begin{enumerate}
              \item \lstasm|MOV AX, [SI+DI]|
              \item \lstasm|MOV AX, [BX+DI]|
              \item \lstasm|MOV AX, [BP+DI]|
              \item \lstasm|MOV AX, [BX+SI]|
          \end{enumerate}

    \item Di antara pasangan instruksi berikut, yang paling cepat
          dieksekusi adalah \dots
          \begin{enumerate}
              \item \lstasm|MOV AX, [BX+SI+2]|
              \item \lstasm|MOV AX, [BP+2]|
              \item \lstasm|MOV AX, [1234h]|
              \item \lstasm|MOV AX, BX|
          \end{enumerate}

    \item Diketahui \reg{BX} = \hex{0100} dan \reg{SI} =
          \hex{0010}. Alamat efektif operand
          \lstasm|MOV AX, [BX+SI+8]| adalah \dots
          \begin{enumerate}
              \item \hex{0110} \item \hex{0118}
              \item \hex{0108} \item \hex{0128}
          \end{enumerate}
\end{enumerate}

\subsection{Soal Uraian}

\begin{enumerate}
    \item Jelaskan mengapa 8086 hanya mengizinkan \reg{BX},
          \reg{BP}, \reg{SI}, dan \reg{DI} sebagai penunjuk alamat
          serta mengapa \code{[AX]} dan \code{[SI+DI]} tidak dapat
          dikodekan. Kaitkan dengan byte \eng{ModR/M} pada
          Bab~\ref{chap:format-instruksi}.

    \item Tuliskan satu operasi yang sama, yaitu menyalin elemen
          berindeks pertama dari sebuah array \dir{DW} ke register
          \reg{AX}, dalam lima mode pengalamatan berbeda. Untuk
          setiap cara, sebutkan segmen \eng{default} yang dipakai.

    \item Diketahui \reg{DS} = \hex{1500}, \reg{SS} = \hex{2500},
          \reg{BX} = \hex{0040}, \reg{BP} = \hex{0060}, dan
          \reg{DI} = \hex{0006}. Hitung alamat efektif dan alamat
          fisik untuk \lstasm|MOV AX, [BX+DI+2]| dan
          \lstasm|MOV AX, [BP]|. Sebutkan pula langkah verifikasi yang
          Anda lakukan pada jendela memori \eng{debugger}.

    \item Tabel nilai mahasiswa berisi \eng{record} enam byte: empat
          byte NIM dan dua byte nilai. Rancang mode pengalamatan untuk
          membaca medan nilai mahasiswa ke-$i$, tuliskan potongan
          kodenya, dan jelaskan mengapa \lstasm|[BX+SI+4]| lebih
          tepat daripada tiga instruksi \lstasm|MOV| terpisah
          (Bab~\ref{chap:prosedur}).
\end{enumerate}

\subsection{Kisi-kisi Penilaian}

\begin{table}[H]
    \centering
    \caption{Kisi-kisi evaluasi Bab~\ref{chap:mode-pengalamatan}.}
    \label{tab:kisi-bab8}
    \small
    \begin{tabular}{@{}p{1.5cm}p{5.4cm}p{3.0cm}c@{}}
        \toprule
        \textbf{Sub-CPMK} & \textbf{Indikator yang diukur} &
        \textbf{Bentuk soal} & \textbf{Bobot} \\
        \midrule
        2.2 & Mengenali mode pengalamatan dari penulisan operand &
        Pilihan ganda 1; Uraian 2 & 25\% \\
        \addlinespace
        2.2 & Menghitung alamat efektif dan alamat fisik &
        Pilihan ganda 5; Uraian 3 & 25\% \\
        \addlinespace
        2.2 & Menentukan segmen \eng{default} dan
        \eng{override} & Pilihan ganda 2; Uraian 3 & 20\% \\
        \addlinespace
        2.2 & Menganalisis legalitas operand dan efisiensi mode &
        Pilihan ganda 3, 4; Uraian 1 & 15\% \\
        \addlinespace
        2.2 & Menerapkan mode pengalamatan pada array dan
        \eng{record} & Uraian 4 & 15\% \\
        \bottomrule
    \end{tabular}
\end{table}

\begin{checklistbox}
\textbf{Checklist Pencapaian Kompetensi Bab~\ref{chap:mode-pengalamatan}}

Seorang mahasiswa dinyatakan kompeten pada Sub-CPMK 2.2 apabila:
\begin{itemize}
    \item[\checked] dapat menyebutkan tujuh mode pengalamatan beserta
          contoh dan kegunaannya;
    \item[\checked] dapat menentukan segmen \eng{default} setiap mode
          serta menuliskan \eng{segment override prefix};
    \item[\checked] dapat menghitung alamat efektif dan alamat fisik
          operand memori tanpa kesalahan;
    \item[\checked] dapat menjelaskan alasan penolakan operand
          \code{[AX]}, \code{[BX+BP]}, dan \code{[SI+DI]};
    \item[\checked] dapat memilih mode yang efisien untuk akses array
          dan \eng{record};
    \item[\checked] dapat memverifikasi perhitungan alamatnya melalui
          jendela memori \eng{debugger}.
\end{itemize}
\end{checklistbox}


\end{document}
